\scanpage{9}
- \mspage{6} -
\[
(1.8)\qquad u^4c'_4=c_4\qquad u^6c'_6=c_6\qquad u^{12}\Delta'=\Delta.
\]
Les formules (1.8) expriment que la section $c_4\pi^{\otimes4}$ (resp. $c_6\pi^{\otimes6}$, \unsure{rosp.|resp.} $\Delta\pi^{\otimes12}$) du fibré $\omega^{\otimes4}$ (resp. $\omega^{\otimes6}$, $\omega^{\unsure{\circ|\otimes}12}$) ne dépend pas du choix arbitraire de $\pi$, $x$ et $y$.

2.\quad \uline{Cas où 2 et 3 sont inversibles}.

Lorsque 2 et 3 sont inversibles, on peut d'une et d'une seule façon choisir $x$ et $y$ ($\pi$ étant donné) de sorte que dans l'équation (1.3.) on ait $a_1=a_2=a_3=0$. On a alors
\[
(2.1)\qquad y^2=x^3+a_4x+a_6,\ \text{avec}\quad\pi=dx/2y,
\]
\[
(2.2)\qquad c_4=-2^4\cdot3a_4,\quad c_6=-2^5\cdot3^3a_6,\quad\Delta=2^4(-4a_4^3-27a_6^2).
\]
Posant $Y=2y$, $X=x$, on retombe sur la forme de Weierstrass
\[
(2.3)\qquad Y^2=4X^3-g_2X-g_3,\quad\text{avec}\quad\pi=dX/Y,\quad\text{et}
\]
\[
(2.4)\qquad g_2=-4a_4=\frac1{12}c_4\qquad g_3=-4a_6=\frac1{216}c_6
\]
\[
\Delta=g_2^3-27g_3^2.
\]
\uline{Proposition} 2.5.\quad \uline{Au dessus de} $\mathrm{Spec}(\mathbb Z[2^{-1},3^{-1}]$, \uline{il existe un schéma de modules pour les courbes de genre un munies d'une forme différentielle invariante inversible}. \uline{\unsure{C;|Ce} schéma est}
\[
\mathrm{Spec}(\mathbb Z[2^{-1},3^{-1}][g_2,g_3],
\]
\margin{8}
